
% BEGIN bib
 
\par\medskip\noindent\textbf{References}\par
  
\par\noindent  [1]  A. Abouelaoualim, K.Ch. Das, L. Faria, Y. Manoussakis, C. Martinhon, and R. Saad.  Paths and trails in edge-colored graphs.  Theoretical Computer Science, 409(3):497–510, 2008.  doi:10.1016/j.tcs.2008.09.021.  
 
\par\noindent  [2]  Adrian Bondy and U. S. R. Murty.  Graph Theory.  Number 1 in Graduate Texts in Mathematics. Springer London, 2008.  URL: https://link.springer.com/book/9781846289699.  
 
\par\noindent  [3]  Jørgen Bang-Jensen and Gregory Z Gutin.  Digraphs: theory, algorithms and applications.  Springer Science \& Business Media, 2008.  
 
\par\noindent  [4]  Claude Berge.  Two theorems in graph theory.  Proceedings of the National Academy of Sciences of the United States of America, 43(9):842–844, September 1957.  URL: https://www.jstor.org/stable/89875.  
 
\par\noindent  [5]  Pierre-Louis Curien.  Introduction to linear logic and ludics, part II.  lecture notes cs/0501039, arXiv, 2005.  URL: https://arxiv.org/abs/cs/0501039.  
 
\par\noindent  [6]  Vincent Danos.  La Logique Linéaire appliquée à l’étude de divers processus de normalisation (principalement du $\lambda$-calcul).  Thèse de doctorat, Université Paris VII, 1990.  
 
\par\noindent  [7]  Vincent Danos and Laurent Regnier.  The structure of multiplicatives.  Archive for Mathematical Logic, 28:181–203, 1989.  doi:10.1007/BF01622878.  
 
\par\noindent  [8]  Paulin Jacobé de Naurois and Virgile Mogbil.  Correctness of linear logic proof structures is NL-complete.  Theoretical Computer Science, 412(20):1941–1957, 2011.  doi:10.1016/J.TCS.2010.12.020.  
 
\par\noindent  [9]  Rémi Di Guardia.  Identity of Proofs and Formulas using Proof-Nets in Multiplicative-Additive Linear Logic.  Thèse de doctorat, École Normale Supérieure de Lyon, September 2024.  
 
\par\noindent  [10]  Thomas Ehrhard.  A new correctness criterion for MLL proof nets.  In Thomas A. Henzinger and Dale Miller, editors, Joint Meeting of the Twenty-Third EACSL Annual Conference on Computer Science Logic (CSL) and the Twenty-Ninth Annual ACM/IEEE Symposium on Logic in Computer Science (LICS), pages 38:1–38:10. ACM, 2014.  doi:10.1145/2603088.2603125.  
 
\par\noindent  [11]  Arnaud Fleury and Christian Retoré.  The mix rule.  Mathematical Structures in Computer Science, 4(2):273–285, 1994.  doi:10.1017/S0960129500000451.  
 
\par\noindent  [12]  Shinya Fujita, Ruonan Li, and Shenggui Zhang.  Color degree and monochromatic degree conditions for short properly colored cycles in edge-colored graphs.  Journal of Graph Theory, 87(3):362–373, 2018.  doi:10.1002/jgt.22163.  
 
\par\noindent  [13]  Hortensia Galeana-Sánchez, Rocío Rojas-Monroy, Rocío Sánchez-López, and Juana Imelda Villareal-Valdés.  H-cycles in H-colored multigraphs.  Graphs and Combinatorics, 38:62, 2022.  doi:10.1007/s00373-022-02464-4.  
 
\par\noindent  [14]  Jean-Yves Girard.  Linear logic.  Theoretical Computer Science, 50:1–102, 1987.  doi:10.1016/0304-3975(87)90045-4.  
 
\par\noindent  [15]  Jean-Yves Girard.  Geometry of interaction I: an interpretation of system ${F}$.  In Ferro, Bonotto, Valentini, and Zanardo, editors, Logic Colloquium ’88. North-Holland, 1988.  
 
\par\noindent  [16]  Jean-Yves Girard.  Quantifiers in linear logic II.  In Corsi and Sambin, editors, Nuovi problemi della logica e della filosofia della scienza, pages 79–90, Bologna, 1991. CLUEB.  
 
\par\noindent  [17]  Jean-Yves Girard.  Linear logic: its syntax and semantics.  In Jean-Yves Girard, Yves Lafont, and Laurent Regnier, editors, Advances in Linear Logic, volume 222 of London Mathematical Society Lecture Note Series, pages 1–42. Cambridge University Press, 1995.  
 
\par\noindent  [18]  Jean-Yves Girard.  Proof-nets: the parallel syntax for proof-theory.  In Aldo Ursini and Paolo Agliano, editors, Logic and Algebra, volume 180 of Lecture Notes In Pure and Applied Mathematics, pages 97–124, New York, 1996. Marcel Dekker.  
 
\par\noindent  [19]  Jerrold W. Grossman and Roland Häggkvist.  Alternating cycles in edge-partitioned graphs.  Journal of Combinatorial Theory, Series B, 34(1):77–81, 1983.  doi:10.1016/0095-8956(83)90008-4.  
 
\par\noindent  [20]  Giulio Guerrieri, Giulia Manara, Lorenzo Tortora de Falco, and Lionel Vaux Auclair.  Confluence for proof-nets via parallel cut elimination.  In Nikolaj S. Bjørner, Marijn Heule, and Andrei Voronkov, editors, LPAR 2024: Proceedings of 25th Conference on Logic for Programming, Artificial Intelligence and Reasoning, Port Louis, Mauritius, May 26-31, 2024, volume 100 of EPiC Series in Computing, pages 464–483. EasyChair, 2024.  doi:10.29007/VKFN.  
 
\par\noindent  [21]  Stefano Guerrini.  Correctness of multiplicative proof nets is linear.  In Proceedings of the fourteenth annual symposium on Logic In Computer Science, pages 454–463, Trento, July 1999. IEEE, IEEE Computer Society Press.  doi:10.1109/LICS.1999.782640.  
 
\par\noindent  [22]  Dominic Hughes.  Simple multiplicative proof nets with units, 2005.  arXiv:math/0507003.  
 
\par\noindent  [23]  Dominic Hughes.  Simple free star-autonomous categories and full coherence.  Journal of Pure and Applied Algebra, 216(11):2386–2410, 2012.  doi:10.1016/j.jpaa.2012.03.020.  
 
\par\noindent  [24]  Dominic Hughes and Rob van Glabbeek.  Proof nets for unit-free multiplicative-additive linear logic.  ACM Transactions on Computational Logic, 6(4):784–842, 2005.  doi:10.1145/1094622.1094629.  
 
\par\noindent  [25]  Anton Kotzig.  On the theory of finite graphs with a linear factor II.  Matematicko-Fyzikálny Časopis, 09(3):136–159, 1959.  In Slovak, with as original title Z teórie konečných grafov s lineárnym faktorom II.  URL: https://eudml.org/doc/29908.  
 
\par\noindent  [26]  Yves Lafont.  From proof nets to interaction nets.  In Jean-Yves Girard, Yves Lafont, and Laurent Regnier, editors, Advances in Linear Logic, volume 222 of London Mathematical Society Lecture Note Series, pages 225–247. Cambridge University Press, 1995.  
 
\par\noindent  [27]  Olivier Laurent.  Polynomial time in untyped elementary linear logic.  Theoretical Computer Science, 813:117–142, 2020.  doi:10.1016/J.TCS.2019.10.002.  
 
\par\noindent  [28]  Ruonan Li and Bo Ning.  A revisit to Bang-Jensen-Gutin conjecture and Yeo’s theorem.  preprint, 2022.  doi:10.48550/arXiv.2207.03793.  
 
\par\noindent  [29]  Paul-André Melliès.  A topological correctness criterion for multiplicative non-commutative logic.  In Thomas Ehrhard, Jean-Yves Girard, Paul Ruet, and Philip J. Scott, editors, Linear Logic in Computer Science, volume 316 of London Mathematical Society Lecture Note Series, pages 283–322. Cambridge University Press, November 2004.  
 
\par\noindent  [30]  François Métayer.  Homology of proof-nets.  Archive for Mathematical Logic, 33(3):169–188, 1994.  doi:10.1007/bf01203031.  
 
\par\noindent  [31]  Gleb Nenashev.  A short proof of Kotzig’s theorem.  preprint, 2014.  doi:10.48550/arXiv.1402.0949.  
 
\par\noindent  [32]  Lê Thành Dũng Nguy$\tilde{\hat{\text{e}}}$n.  Unique perfect matchings, forbidden transitions and proof nets for linear logic with mix.  Logical Methods in Computer Science, 16(1), February 2020.  doi:10.23638/LMCS-16(1:27)2020.  
 
\par\noindent  [33]  Michele Pagani and Lorenzo Tortora de Falco.  Strong normalization property for second order linear logic.  Theoretical Computer Science, 411(2):410–444, 2010.  doi:10.1016/J.TCS.2009.07.053.  
 
\par\noindent  [34]  Christian Retoré.  Handsome proof-nets: perfect matchings and cographs.  Theoretical Computer Science, 294(3):473–488, 2003.  doi:10.1016/S0304-3975(01)00175-X.  
 
\par\noindent  [35]  Paul D. Seymour.  Sums of circuits.  Graph Theory and Related Topics, pages 341–355, 1978.  
 
\par\noindent  [36]  D. J. Shoesmith and T. J. Smiley.  Multiple-Conclusion Logic.  Cambridge University Press, 1978.  
 
\par\noindent  [37]  D. J. Shoesmith and T. J. Smiley.  Theorem on directed graphs, applicable to logic.  Journal of Graph Theory, 3(4):401–406, 1979.  doi:10.1002/jgt.3190030412.  
 
\par\noindent  [38]  Stefan Szeider.  On theorems equivalent with Kotzig’s result on graphs with unique 1-factors.  Ars Combinatoria, 73:53–64, 2004.  URL: https://www.ac.tuwien.ac.at/files/pub/szeider-AC-2004.pdf.  
 
\par\noindent  [39]  Anders Yeo.  A note on alternating cycles in edge-coloured graphs.  Journal of Combinatorial Theory, Series B, 69(2):222–225, 1997.  doi:10.1006/jctb.1997.1728.  
  
% END bib
